% Kapitel 5 -- Von Kuonrats Dienst am Hofe von Hardegg

\section{Von Kuonrats Dienst am Hofe von Hardegg}

\mylettrine{D}{er} junge \gls{Kuonrat von Steinare} weilte nun mehrere Monde am Hofe des Grafen \gls{Kuonrat von Hardegg}. Die Burg \gls{Hardegg} war ihm bald so vertraut wie das väterliche Ottenstein, und doch fremd in einer Weise, die er nicht hätte benennen können: hier schien alles größer, lauter und gieriger. Die Fackeln im großen Saal brannten bis tief in die Nacht, das Fleisch auf den Tischen war reichlicher, der Wein herber, und der Lärm des Gesindes vermischte sich am Abend mit den Windstoßen, die durch die Thayaschlucht fuhren und die Torflügel in den Angeln ächzen ließen.

Was dem Knaben in den frühen Wochen als erstes auffiel, war die Haltung der Männer des Hofes gegenüber dem Grafen. Knechte und Bewaffnete zeigten ihre Unterwerfung nach Art der Männer: mit straffer Schulter, kurzem Nicken, gehaltenen Schritten. Ihre Augen begegneten jenen des Grafen manchmal, dann wichen sie aus; sie sprachen knapp und mit Bedacht, wählten ihre Worte wie einer, der auf gefrorenem Teich geht und jeden Schritt prüft. Der Knabe erkannte in alledem das Wesen des Gehorsams, wie ihn sein Vater \gls{Hadmar von Steinare} ihn beschrieben hatte: ein Dienst, der dem Herrn gebührt, erlangt durch Verdienst und Stand.

Aber die Frauen des Hofes verhielten sich anders. Und hier, an diesem Punkt, muss ich als Chronist dieser Geschichte innehalten, denn hier liegt jener Irrtum, der sich tiefer in \gls{Kuonrat von Steinare}s Herz eingrub als manch spätere Wunde.

Die Hofdamen und Mägde, die im Dienst auf \gls{Hardegg} standen, zeigten dem Grafen eine Unterwerfung, die von jener der Männer grundverschieden war, obwohl beide denselben Ursprung hatten. Wenn der Graf den Saal betrat, wurden sie still; ihr Geplauder erstarb wie eine Kerzenflamme im Wind. Manche senkten die Augen sogleich und hielten sie gesenkt, als wäre ihr Blick zu schwer für ihn. Andere röteten sich leicht, wenn er vorüberging, ohne sie eines Blickes zu würdigen. Die Schritte der Mägde wurden hastig, wenn der Graf nach etwas verlangte; sie beeiferten sich, ihm zu gefallen, und näherten sich ihm mit weichen, beruhigenden Stimmen, wie man mit einem Hund spricht, dessen Laune unbekannt ist.

\gls{Kuonrat von Steinare} sah all dies mit den Augen eines Vierzehnjährigen, dessen Herz bereits weit offenstand für das Geschlecht, das er über alles Irdische stellte. Von Frauen hatte er sich eine Vorstellung gebildet, die mit dem wirklichen Weibe wenig gemein hatte, wohl aber mit den Liedern, die man ihm auf Ottenstein vorgesungen hatte. Für ihn war das Weib ein hohes Wesen, strahlend und unergründlich, eine Art irdische Gottheit, vor der ein Mann sich bewähren musste, um ihrer Gunst würdig zu werden. Er dachte an Frauen oft und heimlich, und dieser Gedanke war für ihn keine Sünde, sondern der süßeste Antrieb seines Lebens.

Nun sah er die Frauen des Hofes erblassen und erröten und verstummen vor dem Grafen, und er las es falsch, wie ein Kind liest, dem die Buchstaben noch fremd sind. Er dachte, die Frauen verneigten sich vor dem Grafen aus demselben Grund, aus dem eine Blume sich der Sonne zuneigt: weil sie seine Macht spürten und in ihr das Begehrenswerte erkannten. Das Erröten der Mägde deutete er als stilles Sehnen; das Senken der Augen als Bescheidenheit vor einem Manne, der ihnen zu würdig war für ihren Blick; das eifrige Bedienen als den Wunsch, seine Gunst zu erwerben. Die Blässe, die manche zeigten, wenn er sie ansprach, hielt er für jene Erschütterung, von der die Lieder singen: das Herzklopfen des Weibes, das den Rechten erblickt.

So heißt es, und ich glaube es gern, denn ich kenne genug von dieser Welt, um zu wissen: ein Knabe, der Frauen ersehnt, sieht Begehren, wo er hinsieht. Und er sah es überall.

Es gab einen Abend, so ward mir berichtet, an dem der Irrtum besonders tief in den Knaben eingedrungen sein muss. Der Graf \gls{Kuonrat von Hardegg} hielt Tafel, und Ritter aus den umliegenden Gütern waren geladen, manche mitsamt ihren Frauen. Der Saal war voll und warm vom Feuer; Fett tropfte vom Braten, der Wein kreiste. \gls{Kuonrat von Steinare} saß, wie es ihm der Graf gestattete, nahe genug am hohen Tisch, um alles sehen und hören zu können, und er saß dort mit der stillen Würde eines Knaben, dem nicht aufgeht, dass er aus Gunst sitzt und nicht aus Recht.

Unter den geladenen Frauen befand sich die junge Ehefrau eines Ritters aus der Gegend, eine stille, blasse und sehr dünne Person, die, so heißt es, für ihre Schönheit bekannt war und für ihr zurückhaltendes Wesen. Sie saß nahe dem Grafentisch, die Hände im Schoß gefaltet, der Blick kaum wandernd. Als der Graf \gls{Kuonrat von Hardegg} einmal zu ihr hinübersah und eine Bemerkung über den Wein machte, den sie halte, wurde sie sichtlich blasser und lächelte mit dem Mund, nicht mit den Augen.

Dies sah \gls{Kuonrat von Steinare}, und er sah es mit seiner eigentümlichen Blindheit. Er sah die Blässe als Erregung; er sah das Lächeln als die Verborgenheit weiblicher Zuneigung, jene Keuschheit, die das Begehren hinter Bescheidenheit versteckt. Er sah die Stille, mit der sie dasaß, als das Schweigen einer Frau vor einem Manne, vor dem die Worte versagen: nicht weil er Furcht einflößt, sondern weil er die Seele bewegt.

Der Ehemann der Frau hatte bemerkt, dass der Graf sie angesehen hatte; er sprach danach lauter und aß hastiger, wie einer, der Angst bekämpft durch Geräusch. Aber auch dies entging dem Knaben, oder er deutete es als die bäurische Eifersucht des Schwachen, die ihm nur bewies, dass der Graf mehr galt als der Ritter, und dass Frauen dies wussten.

Gegen Ende des Abends ließ \gls{Kuonrat von Hardegg} eine Magd herankommen, die ihm den Becher füllen sollte. Das Mädchen, kaum älter als der Knabe selbst, hielt den Krug mit beiden Händen, und doch zitterten diese Hände so, dass der Wein leise gegen das Zinn schlug. Als sie fertig war, wollte sie schnell zurückweichen; der Graf aber fasste sie einen Moment fest, ohne Güte, ohne Wärme, nur als Zeichen, dass er es konnte und aus reiner Lüsternheit. Das Mädchen erstarrte wie ein Tier, das einen Falken über sich wähnt.

\gls{Kuonrat von Steinare} sah dies, und er dachte, das Zittern des Mädchens sei Aufregung über die Berührung des Herrn. Er dachte, das Erstarren sei Entzücken, jene Lähmung, die den Menschen befällt, wenn ein großes Glück ihn überwältigt. Er hatte in den Liedern gehört, dass Frauen beim Anblick des Geliebten die Glieder schwer werden und die Seele aufsteigt; hier sah er es, meinte er, mit eigenen Augen.

Der Herr allein weiß, was in dem Knaben in jener Stunde wirklich vorging. Doch was aus diesem Abend erwuchs, lässt sich in seinen späteren Jahren ablesen.

Der Graf \gls{Kuonrat von Hardegg} sprach manchmal zu dem Knaben über Frauen, denn er mochte lehren, sofern das Lehren ihm Gelegenheit gab, seine eigene Urteilskraft zu preisen. Diese Lehrsprüche hatten stets den Ton des Verachtenden: das Weib sei ein nützliches Werkzeug, weich dort, wo es nützlich war, und gefügig durch den richtigen Druck. \enquote{Wer die Weiber kennt}, sprach er einmal zu dem Knaben, während er seinen Becher drehte, \enquote{der weiß, dass sie Worte verabscheuen und Taten achten. Zeige Stärke, und sie beugen sich. Zeige Weichheit, und sie verlachen dich.}

Der Knabe hörte dies und verstand es mit jener besonderen Fähigkeit des Irrtümlichen, Worte so zu hören, wie man sie hören will. Er verstand: zeige Stärke, und Frauen begehren dich. Er verstand nicht: zeige Stärke, und Frauen fürchten dich zu sehr, um zu widersprechen. Beides führt zu Beugen und Schweigen; nur das Innere ist verschieden, und das Innere sieht man nicht, wenn man es nicht sehen will.

Was den Knaben am nachhaltigsten formte, war denn auch nicht das, was der Graf sagte, sondern was er war: ein Mann, vor dem die Männer ihre Stimmen senkten und die Frauen erblassten. Und da der Knabe die Unterwerfung der Männer kannte und richtig las, aber das Erblassen der Frauen falsch deutete, wurde ihm das Erblassen wichtiger. Denn er selbst sehnte sich nicht nach der stummen Achtung der Männer. Er sehnte sich, von Grund auf und mit jeder Faser seines törichten jungen Herzens, nach der Zuneigung der Frauen. Das war, so glaube ich gern, die einzige wirklich lebendige Sehnsucht in ihm, lauter als all sein Hochmut.

So wuchs in dem Jungen ein Bild: ein Herr, der Frauen bezaubert, tut es durch Kraft und Unnahbarkeit. Und er beschloss, dieses Bild in sich aufzunehmen, so gut er konnte. Er begann, den Grafen nachzuahmen in Haltung und Schritt, in der Art, wie man einen Becher hält, in dem Schweigen, das er für Würde hielt. Ob die Mägde, die ihm danach den Blick senkten, darin Herrschaft sahen oder nur einen Knaben, der zu groß zu gehen versuchte, das ist mir nicht überliefert worden.

Der Graf \gls{Kuonrat von Hardegg} behandelte \gls{Kuonrat von Steinare} in diesen Monaten mit jener kalten Freundlichkeit, die ein erfahrener Jäger einem jungen Hund entgegenbringt: er zeigte ihm Gutes und Strenge im Wechsel, lobte ihn vor anderen und tadelte ihn danach allein, ließ ihm manches durchgehen, was er einem anderen nicht gegönnt hätte, und prüfte dabei stets, wie tief die Unterwerfung reichte. \enquote{Du bist jung und eitel, Kuonr\^{a}t}, sprach er einmal zu ihm, \enquote{und du glaubst, dass dir die Welt aus Liebe gehorcht. Das ist die süßeste aller Torheiten.} Doch die Nähe des Herrn tränkte den Knaben wie Wein, und er sog begierig jeden Lobpreis auf und ließ die Warnung vorbeiziehen wie Rauch.

Am Ende der Tage, wenn die Launen des Grafen wechselten, zeigte sich dem Knaben, was der Hof in lichteren Stunden verbarg: die Kälte, die Strenge, die Bereitschaft zur Strafe. Einmal ließ der Graf einen Diener vor versammeltem Gefolge züchtigen, weil dieser einen Trunk verschüttet hatte; und das Mädchen, das danach den Becher neu füllte, bewegte sich so lautlos durch den Saal, als wolle sie unsichtbar werden. \gls{Kuonrat von Steinare} sah auch dies. Und er sah nicht Schrecken in dem Mädchen; er sah Ehrerbietung. Er sah nicht das Unsichtbarwerden-Wollen der Angst, er sah die Bescheidenheit der Bewunderung.

Noch fehlte ihm jeder Sinn dafür, dass ein erbleichendes Mädchen, eine schweigende Frau, eine stille und flinke Dienerin vielleicht nichts anderes hegte als den Wunsch, dieser Stunde unbeschadet zu entkommen. Er sah die Frucht und kannte den Baum nicht.

So setzte sich die Saat des Irrtums in seinem Herzen: er glaubte, die Frauen verneigten sich vor dem Grafen, weil sie ihn begehrten, und er glaubte, er selbst müsse nur werden wie der Graf, dann würden auch sie sich vor ihm neigen aus freiem Willen und in Zuneigung. Was für ein finsteres Geschenk ihm der Hof von \gls{Hardegg} damit mitgab, sollte sich erst mit der Zeit zeigen.
